\section*{Hoofdstuk \ref{ch:b1:genomes} --- Genomen van cellen en virussen}

\begin{solution}{exo:b1:genomes:1}
Bacterie: één ringvormig \omterm{def:b1:genomes:genome}{chromosoom} (plus \omterm{def:b1:genomes:genome}{plasmiden}) in een \omterm{def:b1:genomes:genome}{nucleoïde} zonder
eigen membraan, samengepakt door superspiralisatie en kleine basische \omterm{def:b1:proteins:peptide}{eiwitten}, voor
ongeveer \qty{90}{\%} coderend. Eukaryoot: verscheidene lineaire \omterm{def:b1:genomes:genome}{chromosomen} in een
celkern, om \omterm{def:b1:genomes:nucleosome}{histonen} gewonden tot \omterm{def:b1:genomes:nucleosome}{chromatine}, voor enkele procenten coderend.
\end{solution}

\begin{solution}{exo:b1:genomes:2}
147 \omterm{thm:b1:nucleic-acids:helix}{basenparen} 1,7 keer om een achttal \omterm{def:b1:genomes:nucleosome}{histonen} gewonden (twee elk van H2A, H2B,
H3, H4), door H1 afgesloten, om de 200 paren één. Kralen aan een snoer $\times 6$;
vezel van \qty{30}{nm} $\times 40$; lussen $\times 1000$; metafasechromosoom
$\times \num{10000}$.
\end{solution}

\begin{solution}{exo:b1:genomes:3}
\qty{1.5}{\%}; transponeerbare elementen en hun resten, \qty{45}{\%}.
\end{solution}

\begin{solution}{exo:b1:genomes:4}
Een \omterm{def:b1:genomes:genome}{genoom} (\omterm{def:b1:nucleic-acids:chain}{DNA} of \omterm{def:b1:nucleic-acids:chain}{RNA}) in een eiwitcapside, soms met envelop, dat zich uitsluitend
binnen een \omterm{def:b1:cell-unit-of-life:cell}{cel} voortplant met de machinerie van die \omterm{def:b1:cell-unit-of-life:cell}{cel}. Buiten een \omterm{def:b1:cell-unit-of-life:cell}{cel} heeft het
geen stofwisseling, geen groei, geen reactie --- een levenloos deeltje; het
``leeft'' alleen doordat het erfelijke informatie draagt die evolueert.
\end{solution}

\begin{solution}{exo:b1:genomes:5}
$\num{4.6e6}\times 0.34 = \qty{1.56}{mm}$; $\num{4.6e6}/200 = \num{23000}$
nucleosomen. In plaats daarvan: negatieve superspiralisatie door het gyrase en
opvouwing in zo'n vijftig lussen op een eiwitkern, met kleine basische \omterm{def:b1:proteins:peptide}{eiwitten} die
het \omterm{def:b1:nucleic-acids:chain}{DNA} buigen --- duizendvoudig, tot een \omterm{def:b1:genomes:genome}{nucleoïde} van een micrometer.
\end{solution}

\begin{solution}{exo:b1:genomes:6}
\omterm{def:b1:genomes:gene}{Exonen} $1500/8 = \qty{190}{bp}$; \omterm{def:b1:genomes:gene}{intronen} $25\,500/7 = \qty{3640}{bp}$;
$25.5/27 = \qty{94}{\%}$ van het transcript wordt weggegooid.
\end{solution}

\begin{solution}{exo:b1:genomes:7}
\emph{E. coli} 930 \omterm{def:b1:genomes:gene}{genen} per Mb; gist 500; vlieg 100; mens 6; longvis 0,15. De
gendichtheid daalt met vier ordes van grootte terwijl het aantal \omterm{def:b1:genomes:gene}{genen} vijfvoudig
stijgt: het extra \omterm{def:b1:nucleic-acids:chain}{DNA} is niet-genisch.
\end{solution}

\begin{solution}{exo:b1:genomes:8}
$10^6\times 300 = \qty{3e8}{bp}$: \qty{9}{\%} van het \omterm{def:b1:genomes:genome}{genoom}. $60\times
10^6/20 = \num{3e6}$ generaties: gemiddeld één nieuwe invoeging per drie
generaties, ergens in de lijn.
\end{solution}

\begin{solution}{exo:b1:genomes:9}
De eicel levert het \omterm{def:b1:cell-unit-of-life:cell}{cytoplasma} met haar \omterm{def:b1:eukaryotic-cell:mitochondrion}{mitochondriën}; de weinige \omterm{def:b1:eukaryotic-cell:mitochondrion}{mitochondriën} van
de zaadcel worden buitengesloten of na de bevruchting vernietigd. Het
mitochondriale \omterm{def:b1:nucleic-acids:chain}{DNA} gaat daardoor ongemengd van moeder op kind, en de mutaties die
zich erin ophopen tekenen de moederlijke lijnen door de tijd heen.
\end{solution}

\begin{solution}{exo:b1:genomes:10}
Gevoede kweek, twee uur: lysogenie geeft $2^6 = 64$ kopieën (één per gastheerdeling);
lysis geeft 100 fagen in \qty{40}{min}, en hun nakomelingen $100^3 = 10^6$ in twee
uur als er gastheren in overvloed zijn --- lysis wint met ruime afstand. Hongerende
kweek: lysogenie geeft één kopie, veilig binnen een levende \omterm{def:b1:cell-unit-of-life:cell}{cel}; lysis geeft 100
fagen zonder \omterm{def:b1:cell-unit-of-life:cell}{cel} om te besmetten, die vergaan. De schakelaar leest de toestand van
de gastheer omdat de beste strategie ervan afhangt of er nieuwe gastheren
beschikbaar zullen zijn.
\end{solution}

\begin{solution}{exo:b1:genomes:11}
$\num{1.3e11}\times 0.34\,\mathrm{nm} = \qty{44}{m}$ (haploïd);
$\num{6.5e8}$ nucleosomen; \omterm{def:b1:genomes:nucleosome}{histonen} $\num{6.5e8}\times 108\,000\times
1.66\times 10^{-24} = \qty{117}{pg}$. Oorsprongen: $\num{1.3e6}$, die elk aan
weerszijden \qty{50}{kb} repliceren bij \qty{50}{bp/s}: \qty{1000}{s} --- de tijd is
bij genoeg oorsprongen niet het probleem; de kosten zijn de \omterm{def:b1:nucleic-acids:nucleotide}{nucleotiden}
(\qty{130}{pg} \omterm{def:b1:nucleic-acids:chain}{DNA} per deling), de \omterm{def:b1:genomes:nucleosome}{histonen}, de tijd om ze te maken, en een celkern
en een \omterm{def:b1:cell-unit-of-life:cell}{cel} die vele malen groter zijn, en daarom delen zulke \omterm{def:b1:cell-unit-of-life:cell}{cellen} traag.
\end{solution}

\begin{solution}{exo:b1:genomes:12}
Een ontwerp zou bevatten wat nodig is en niets meer; een verslag bevat wat er is
gebeurd. Pseudogenen zijn de lijken van verdubbelde \omterm{def:b1:genomes:gene}{genen} die door mutatie zijn
gestorven; transponeerbare elementen zijn parasieten die zichzelf tientallen
miljoenen jaren hebben gekopieerd en nooit zijn verwijderd; genfamilies tonen
verdubbeling gevolgd door uiteenlopen; en de C-waardeparadox toont dat de
hoeveelheid \omterm{def:b1:nucleic-acids:chain}{DNA} de geschiedenis van ophoping en verlies in een lijn weerspiegelt en
niet haar behoeften. De selectie houdt de \omterm{def:b1:genomes:gene}{genen} werkend en laat de rest driften:
wat wij in een \omterm{def:b1:genomes:genome}{genoom} lezen is grotendeels geschiedenis.
\end{solution}

\begin{solution}{pb:b1:genomes:1}
\textbf{1.} $\num{6.4e9}\times 0.34\,\mathrm{nm} = \qty{2.18}{m}$.
\textbf{2.} $\num{6.4e9}/10 = \num{6.4e8}$ windingen.
\textbf{3.} $\num{2.5e8}\times 0.34 = \qty{8.5}{cm}$; \qty{1.7}{cm}.
\textbf{4.} $\frac{4}{3}\pi\times 3^3 = \qty{113}{\micro m^3}$.
\textbf{5.} $\pi\times(1\,\mathrm{nm})^2\times 2.18\,\mathrm{m} =
\qty{6.8e-18}{m^3} = \qty{6.8}{\micro m^3}$: \qty{6}{\%} van de celkern.
\textbf{6.} $\num{6.4e9}\times 650\times 1.66\times 10^{-24} =
\qty{6.9e-12}{g} = \qty{6.9}{pg}$.
\textbf{7.} Kernmassa $113\times 10^{-12}\,\mathrm{mL}\times 1.1 =
\qty{124}{pg}$; \omterm{def:b1:nucleic-acids:chain}{DNA} is er \qty{5.6}{\%} van, \omterm{def:b1:proteins:peptide}{eiwit} \qty{20}{\%}.
\textbf{8.} $\num{6.4e9}/200 = \num{3.2e7}$.
\textbf{9.} $\num{3.2e7}\times 108\,000\times 1.66\times 10^{-24} =
\qty{5.7}{pg}$: ongeveer evenveel als het \omterm{def:b1:nucleic-acids:chain}{DNA}.
\textbf{10.} $68/11 = 6.2$.
\textbf{11.} Zes nucleosomen, \qty{1200}{bp} $= \qty{408}{nm}$ helix, in
\qty{11}{nm}: factor $37$.
\textbf{12.} $2.18/37 = \qty{5.9}{cm}$ vezel --- tienduizend keer de
kerndoorsnede, daarbinnen opgevouwen.
\textbf{13.} Alle $\num{3.2e7}$, plus evenveel nieuwe op de tweede kopie:
$\num{6.4e7}$ in \qty{28800}{s}, ongeveer \num{2200} per seconde.
\textbf{14.} $75\,000\times 0.34/37 = \qty{690}{nm}$ vezel per lus;
$\num{6.4e9}/75\,000 = \num{85000}$ lussen.
\textbf{15.} $\qty{8.5}{cm}/\qty{10}{\micro m} = 8500$.
\textbf{16.} $\pi\times 0.35^2\times 10 = \qty{3.85}{\micro m^3}$; \omterm{def:b1:nucleic-acids:chain}{DNA}
$\pi\times 10^{-6}\times 0.085\,\mathrm{m} = \qty{0.27}{\micro m^3}$:
\qty{7}{\%} --- de rest is histon- en \omterm{def:b1:proteins:peptide}{geraamte-eiwit} en water.
\textbf{17.} Totaal \omterm{def:b1:nucleic-acids:chain}{DNA} $\num{6.4e9}$ paren bij dezelfde dichtheid
($\num{2.5e8}$ paren in \qty{3.85}{\micro m^3}, per chromatide; twee chromatiden
per \omterm{def:b1:genomes:genome}{chromosoom}): $2\times 6.4\times 10^9/2.5\times 10^8\times
3.85 = \qty{197}{\micro m^3}$, groter dan de interfasekern omdat de \omterm{def:b1:genomes:genome}{chromosomen} nu
staafjes zijn met ruimte ertussen.
\textbf{18.} \omterm{def:b1:genomes:nucleosome}{Chromatine} in de interfase moet worden gelezen en gekopieerd: de
polymerasen en hun factoren hebben toegang tot het \omterm{def:b1:nucleic-acids:chain}{DNA} nodig, zodat het grootste
deel ervan als de open vezel van \qty{30}{nm} en als lussen blijft; \omterm{def:b1:genomes:nucleosome}{chromatine} in
de metafase is inert en kan voor het vervoer strak worden gepakt.
\textbf{19.} $\num{6.4e9}$ plaatsen bij $10^3$ per seconde: $\num{6.4e6}$
seconden, 74 dagen. \omterm{def:b1:proteins:peptide}{Eiwitten} zoeken niet willekeurig: zij binden niet-specifiek aan
het \omterm{def:b1:nucleic-acids:chain}{DNA} en glijden erlangs, en vele kopieën zoeken tegelijk, zodat een plaats
binnen seconden wordt gevonden.
\textbf{20.} $\num{20000}\times 27\,\mathrm{kb} = \qty{540}{Mb}$, \qty{17}{\%} van
\qty{3.2}{Gb}; coderend $\num{20000}\times 1.5 =
\qty{30}{Mb}$, \qty{0.9}{\%} (ongeveer 1,5\% als de kortere \omterm{def:b1:genomes:gene}{genen} zorgvuldiger
worden meegeteld).
\textbf{21.} $113\times 2\times 250/6400 = \qty{8.8}{\micro m^3}$ (twee kopieën);
$\num{5e8}/8.8 = \qty{5.7e7}{bp/\micro m^3}$, hetzelfde als voor de hele celkern.
\textbf{22.} $\num{2.6e11}/\num{5.7e7} = \qty{4600}{\micro m^3}$: een doorsnede van
\qty{21}{\micro m}, in volume veertig keer de menselijke celkern.
\textbf{23.} $\num{3.2e7}$ achttallen, $\num{3.2e10}$ residuen --- een tiende van
de dagelijkse eiwitsynthese van de \omterm{def:b1:cell-unit-of-life:cell}{cel}, besteed aan verpakking.
\textbf{24.} Het \omterm{def:b1:genomes:genome}{chromosoom} van een bacterie is duizend keer korter en
superspiralisatie in lussen volstaat om het in een micrometer te laten passen; een
eukaryoot moet een lengte van tienduizend keer haar celkern opvouwen en heeft een
hiërarchie van verpakking nodig, waarvan het \omterm{def:b1:genomes:nucleosome}{nucleosoom} het eerste niveau is.
\textbf{25.} Ongeveer \num{8500} (\qty{8.5}{cm} in \qty{10}{\micro m}), opgebouwd
als $\times 6$ (nucleosomen), nog eens $\times 6$ (vezel van \qty{30}{nm}, in
totaal $\times 37$), nog eens $\times 27$ (lussen, in totaal $\times 1000$), en nog
eens $\times 8$ (condensatie in de metafase).
\end{solution}
